\section{Risques, arbitrage et feuille de route (G4)}
\label{sec:g4}

\subsection{Scénarios de risque}

Six scénarios traduisent les constats en enjeux métier. La cotation de 1 à 4 suit la notice et le score est le produit de l'impact par la vraisemblance. Le registre complet, avec menaces, faiblesses et mesures existantes, figure en annexe~\ref{ann:risques}.

\begin{figure}[H]
\centering
\begin{minipage}[c]{0.34\linewidth}\centering% Matrice impact × vraisemblance (grille pédagogique 1 à 4 de la notice)
\begin{tikzpicture}[x=9mm,y=7mm,font=\scriptsize]
  \foreach \v in {1,...,4}{\foreach \i in {1,...,4}{
    \pgfmathtruncatemacro\s{\v*\i}
    \ifnum\s>8 \def\c{alerte!22}\else\ifnum\s>4 \def\c{ambre!22}\else\def\c{vert!15}\fi\fi
    \fill[\c] (\v-1,\i-1) rectangle (\v,\i);
    \draw[white, line width=1pt] (\v-1,\i-1) rectangle (\v,\i);
    \node[gris, anchor=north east, font=\tiny] at (\v,\i) {\s};}}
  \newcommand\RisqueCase[3]{\node[font=\scriptsize\bfseries, align=center, text width=8mm] at (#1-0.5,#2-0.5) {#3};}
  \RisqueCase{3}{4}{R01, R02}
\RisqueCase{2}{4}{R03}
\RisqueCase{3}{3}{R04, R06}
\RisqueCase{2}{2}{R05}

  \foreach \v in {1,...,4}{\node at (\v-0.5,-0.3) {\v};}
  \foreach \i in {1,...,4}{\node at (-0.3,\i-0.5) {\i};}
  \node at (2,-0.85) {Vraisemblance};
  \node[rotate=90] at (-0.85,2) {Impact};
\end{tikzpicture}
\end{minipage}\hfill
\begin{minipage}[c]{0.63\linewidth}\small
\begin{tabularx}{\linewidth}{@{}C{0.7cm} Y C{0.35cm} C{0.35cm} C{0.75cm} L{2.3cm}@{}}
\toprule
\entete \hd{ID} & \hd{Scénario} & \hd{I} & \hd{V} & \hd{I×V} & \hd{Décision} \\
\midrule
\lignes{tab_risques_court_sans_actif}
\bottomrule
\end{tabularx}
\end{minipage}
\caption{Matrice impact et vraisemblance, registre résumé (onglet 04).}
\end{figure}

R01 et R02 dominent avec un score de \rscore{R01}. R03 vaut \rscore{R03}, moins que R04 et R06, mais nous le plaçons parmi les risques majeurs pour son impact maximal sur les données de patients et parce qu'il se réduit vite. À score égal, R04 et R06 n'appellent pas le même traitement, le premier relevant de réglages techniques et le second d'une standardisation lourde.

\subsection{Arbitrage sous plafond de \Plafond}

Les montants viennent de \piece{P16} et \piece{T10}. Le lot de restauration n'est compté qu'une fois, car il est commun aux deux pièces. L'outillage quantitatif \piece{(A06)} et la détection outillée \piece{(A11)} sont des alternatives, toutes deux différées, la mesure quantitative relevant du niveau 4.

\begin{figure}[H]
\centering
\resizebox{0.95\linewidth}{!}{% Arbitrage sous plafond (P16, T10, onglet 05), en milliers d'euros
\begin{tikzpicture}[font=\scriptsize, x=0.118cm]
  % échelle
  \foreach \v in {0,20,...,140}{\draw[filet] (\v,0.55) -- (\v,-2.1); \node[text=gris] at (\v,-2.35) {\v};}
  \node[text=gris] at (70,-2.75) {milliers d'euros};
  % plafond
  \draw[alerte, thick] (\PlafondK,0.75) -- (\PlafondK,-2.1); \node[alerte, font=\scriptsize\bfseries, above] at (\PlafondK,0.75) {plafond \PlafondK};
  % retenues
  \node[anchor=east, font=\scriptsize\bfseries] at (-1,0) {Retenu};
  \def\pos{0}
  \foreach \id/\k/\lib [count=\i] in \BudgetRetenu{
    \pgfmathsetmacro\deb{\pos}\pgfmathsetmacro\fin{\pos+\k}
    \pgfmathtruncatemacro\alt{mod(\i,2)}
    \ifnum\alt=1 \def\kcol{vert}\else\def\kcol{vert!70}\fi
    \fill[\kcol] (\deb,-0.35) rectangle (\fin,0.35);
    \node[text=white, font=\scriptsize\bfseries] at ({(\deb+\fin)/2},0.08) {\id};
    \node[text=white, font=\tiny] at ({(\deb+\fin)/2},-0.18) {\k};
    \xdef\pos{\fin}}
  \fill[pattern=north east lines, pattern color=gris] (\pos,-0.35) rectangle (\PlafondK,0.35);
  \draw[gris] (\pos,-0.35) rectangle (\PlafondK,0.35);
  \node[anchor=west, font=\scriptsize] at (121,0) {\BudgetK\ engagés, marge \MargeK};
  % différées
  \node[anchor=east, font=\scriptsize\bfseries] at (-1,-1.35) {Différé};
  \def\pos{0}
  \foreach \id/\k/\lib [count=\i] in \BudgetDiffere{
    \pgfmathsetmacro\deb{\pos}\pgfmathsetmacro\fin{\pos+\k}
    \pgfmathtruncatemacro\alt{mod(\i,2)}
    \ifnum\alt=1 \def\kcol{gris!45}\else\def\kcol{gris!28}\fi
    \fill[\kcol] (\deb,-1.7) rectangle (\fin,-1.0);
    \node[font=\scriptsize\bfseries] at ({(\deb+\fin)/2},-1.27) {\id};
    \node[font=\tiny] at ({(\deb+\fin)/2},-1.53) {\k};
    \xdef\pos{\fin}}
  \node[anchor=west, font=\scriptsize] at (\pos+1,-1.35) {\DiffereK\ reportés en année 2};
\end{tikzpicture}
}
\caption{Actions retenues et différées, en milliers d'euros (onglet 05).}
\end{figure}

\begin{table}[H]
\centering\small
\caption{Plan d'action, synthèse (détail en annexe~\ref{ann:plan})}
\begin{tabularx}{\linewidth}{@{}C{0.9cm} Y R{1.9cm} C{1.4cm} L{3.4cm}@{}}
\toprule
\entete \hd{ID} & \hd{Mesure} & \hd{Coût} & \hd{Choix} & \hd{Pilote} \\
\midrule
\lignes{tab_actions_court}
\midrule
& \textbf{Engagement total} & \textbf{\TotalRetenu} & & Marge \Marge \\
\bottomrule
\end{tabularx}
\end{table}

Le choix traite d'abord R01 et R02 et les processus les plus faibles, comme le recommande la méthode. Chaque report laisse un risque résiduel. Celui de A10 laisse R04 et une part de R02, celui de A04 laisse R06, celui de A11 limite la détection sur R01, celui de A06 prive le pilotage d'indicateurs chiffrés et celui de A07 expose à une faible adhésion. Ces risques doivent être acceptés par écrit par la DG. Les règles FW-31 et FW-17 sont restreintes en régie sans attendre, hypothèse à confirmer avec la DSI. La charge du RSSI à mi-temps reste une fragilité.

\subsection{Feuille de route SSI sur douze mois}

Cette feuille de route concerne MEDICA-SERVICES et non le calendrier de notre projet. Elle suppose une décision du CODIR à la mi-novembre 2026. Les mesures conservatoires sont engagées dès la remise du rapport.

\begin{figure}[H]
\centering
% Feuille de route SSI sur 12 mois (distincte du calendrier de production du projet).
% Hypothèse : décision du CODIR mi-novembre 2026, mois 1 = décembre 2026.
\begin{ganttchart}[
  hgrid={*1{draw=filet, dotted}}, vgrid={*1{draw=filet, dotted}},
  x unit=0.62cm, y unit chart=0.42cm, y unit title=0.45cm,
  title height=1, title label font=\scriptsize\bfseries, title/.append style={fill=accentclair, draw=filet},
  bar/.append style={fill=vert, draw=vert}, bar height=0.55, bar label font=\scriptsize,
]{1}{12}
\gantttitlelist{1,...,12}{1}\\
\newcounter{fr}
\newcommand\FeuilleRoute[4]{\stepcounter{fr}\ganttbar[name=fr\arabic{fr}]{#1}{#2}{#3}%
  \node[anchor=west, font=\tiny, text=gris] at (fr\arabic{fr}.east) {#4};\\}
\FeuilleRoute{Suites de l'investigation}{1}{1}{Clôture de l'incident AL-1048}
\FeuilleRoute{A08 MFA et privilèges}{1}{3}{100 \% MFA ou exception}
\FeuilleRoute{A02 Supervision et reporting}{1}{3}{Alertes tracées}
\FeuilleRoute{A01 Registre des risques}{2}{4}{Registre approuvé}
\FeuilleRoute{A09 Vulnérabilités}{1}{4}{APP01 en M1, SLA en M4}
\FeuilleRoute{A05 Restauration}{3}{6}{9 PV sur 9}
\FeuilleRoute{A03 PSSI et sensibilisation}{4}{6}{PSSI v2 approuvée}
\FeuilleRoute{Stabilisation et mesures}{7}{11}{Indicateurs trimestriels}
\FeuilleRoute{Bilan de maturité et année 2}{12}{12}{Arbitrage A04, A10, A06 ou A11}

\end{ganttchart}

\caption{Feuille de route, mois 1 = décembre 2026.}
\end{figure}

\begin{figure}[H]
\centering
\resizebox{0.9\linewidth}{!}{% Niveaux effectifs démontrés, trajectoire à 12 mois et niveau adéquat (données : scripts/donnees.py)
\begin{tikzpicture}
\begin{axis}[
  xbar, width=0.9\linewidth, height=5.4cm, bar width=7pt,
  xmin=0, xmax=5, xtick={0,1,2,3,4,5}, xlabel={\footnotesize Niveau de maturité},
  ytick={1,...,\MaturiteNb}, yticklabels/.expanded={\MaturiteNoms}, y dir=reverse, ymin=0.5, ymax=\MaturiteNb+0.5,
  yticklabel style={font=\footnotesize}, xticklabel style={font=\footnotesize},
  axis lines*=left, xmajorgrids, grid style={draw=filet},
    legend style={font=\scriptsize, draw=none, at={(0.5,1.02)}, anchor=south, legend columns=-1},
]
\MaturitePlot
\addlegendentry{niveau démontré}
\draw[alerte, thick, dashed] (axis cs:\Cible,0.5) -- (axis cs:\Cible,\MaturiteNb+0.5);
\addlegendimage{alerte, thick, dashed}
\addlegendentry{niveau adéquat (cible \Cible)}
\end{axis}
\end{tikzpicture}
}
\caption{Trajectoire de maturité. Tous les processus au niveau 2 à douze mois, première étape vers la cible.}
\end{figure}
